\documentclass[10pt,svgnames,fragile]{beamer}

% --- PACKAGES --- %
\usepackage[utf8]{inputenc}
\usepackage[T1]{fontenc}
\usepackage[french]{babel}
\usepackage{csquotes}
\usepackage[toc, page]{appendix}
\usepackage{graphicx}
\usepackage{subcaption}
\usepackage{xcolor}
\usepackage{amsmath, amssymb}
\usepackage{amsthm}
\usepackage{mathrsfs}
\usepackage{booktabs}
\usepackage{float}
\usepackage{setspace}
\usepackage{hyperref}
\usepackage{parskip}
\usepackage{caption}
\usepackage{enumitem}
\usepackage{tikz}
\usepackage{tikz-3dplot}
\usepackage{tkz-tab}
\usepackage{dsfont}

% --- MISE EN PAGES --- %
% itemize avec des points
\setlist[itemize,1]{label={\color{black}\small \textbullet}} 
% permet de numéroter les équations en fonction de la section
\numberwithin{equation}{section}
% pour supprimer l'alinéa
\onehalfspacing

% --- QUELQUES COMMANDES --- %
% Ensembles R et N
\newcommand{\bbR}{\mathbb{R}}
\newcommand{\bbC}{\mathbb{C}}
\newcommand{\bbN}{\mathbb{N}}
\newcommand{\bbZ}{\mathbb{Z}}

% Lettres
\newcommand{\calF}{\mathcal{F}}
\newcommand{\calK}{\mathcal{K}}
\newcommand{\calD}{\mathcal{D}}
\newcommand{\calR}{\mathcal{R}}
\newcommand{\calL}{\mathcal{L}}
\newcommand{\calM}{\mathcal{M}}
\newcommand{\calB}{\mathcal{B}}
\newcommand{\calE}{\mathcal{E}}

\newcommand{\scrC}{\mathscr{C}}
\newcommand{\scrS}{\mathscr{S}}
\newcommand{\scrA}{\mathscr{A}}

% Commandes maths
% fonctions vectorielles en gras
\usepackage{bm}
\newcommand{\bfm}[1]{\bm{\mathrm{#1}}}
% divergence
\newcommand{\Div}{\mathrm{div}}
% produit scalaire
\newcommand{\ps}[2]{\langle #1, #2 \rangle}
% norme
\newcommand{\norm}[1]{\left\lVert #1 \right\rVert}
% Commande norme triple
\newcommand{\vertiii}[1]{{\left\vert\kern-0.25ex\left\vert
    \kern-0.25ex\left\vert #1 
\right\vert\kern-0.25ex\right\vert\kern-0.25ex\right\vert}}
% valeur absolue
\newcommand{\abs}[1]{\displaystyle\left\lvert #1 \right\rvert}
% exponentielle
\newcommand{\ex}{\mathrm{e}}
% H^1_0(Omega)
\newcommand{\Hone}[1][\Omega]{H^1_0(#1)}
% D(Omega)
\newcommand{\scrD}{\mathscr{D}(\Omega)}
% opérateurs
\newcommand{\operateurs}{M^{-1}_0 L_0}
% limite
\newcommand{\limm}{\underset{h \rightarrow0}{\lim} \ \dfrac{1}{h}}
% ensemble des fonctions à croissances au plus polynomiales
\newcommand{\scrCinfp}{\scrC^{\infty}_p}
% d droit
\renewcommand{\d}{\mathrm{d}}
% L^1_{loc}
\newcommand{\Lloc}[1]{L^1_{\text{loc}}\big(#1\big)}
% partie réelle
\renewcommand{\Re}{\mathfrak{Re}}
% partie imaginaire
\renewcommand{\Im}{\mathfrak{Im}}
% image
\newcommand{\im}{\mathrm{Im}}

% --- NUMEROTATION THEOREMES --- % 
\newtheorem{theoreme}{Théorème}[section]
\newtheorem{proposition}{Proposition}[section]
\newtheorem{remark}{Remarque}[section]

% --- DEFINITION DE COULEURS --- %
\definecolor{astral}{RGB}{46,116,181}

\usepackage{etex}
\usepackage[normalem]{ulem}
\usepackage{textcomp}
\usepackage{pdflscape}
\usepackage{marvosym}
\usepackage{wasysym}
\usepackage{qtree}
\usepackage{bussproofs}
\usepackage{proof}
\usepackage{fitch}
\usepackage{cancel}
\usepackage{url}
\usepackage{smfthm}


\AtBeginSection[]{
  \begin{frame}
    \footnotesize\tableofcontents[sectionstyle=show/shaded, subsectionstyle=show/show/shaded]
  \end{frame}
}

\AtBeginSubsection[]{
  \begin{frame}
    \footnotesize\tableofcontents[sectionstyle=show/shaded, subsectionstyle=show/shaded/shaded]
  \end{frame}
}

\subtitle{}

\titlegraphic{
\includegraphics[height=2.5cm]{images/logo_um.png} \hspace{1cm} 
\includegraphics[height=2cm]{images/logo_dpt_maths.png}}

\institute[]{} \usetheme{CambridgeUS}
\usepackage{beamer_udl_theme}
\setbeamertemplate{navigation symbols}{%
	%insertslidenavigationsymbol%
	%insertframenavigationsymbol%
	%insertsubsectionnavigationsymbol%
	%insertsectionnavigationsymbol%
	%insertdocnavigationsymbol%
	%insertbackfindforwardnavigationsymbol%
}

\setbeamertemplate{section in toc}{%
  \inserttocsectionnumber.~\inserttocsection
}
\setbeamertemplate{subsection in toc}{%
  \hspace{1.5em}\inserttocsectionnumber.\inserttocsubsectionnumber~\inserttocsubsection\par
}

\usetheme{default}
\author{Quentin LAHAYE}
\date{}
\title{Autour de l'équation de McKendrick-Lotka}

\begin{document}

\maketitle

% ==========================
% Introduction 
% ==========================

\begin{frame}{Introduction}
    \begin{itemize}
        \item L'équation de McKendrick-Lotka est une équation aux dérivées partielles utilisée en dynamique des populations pour modéliser l'évolution d'une population structurée par l'âge. \pause
        \item Une forme classique de cette équation s'écrit
        \begin{equation*}
            \left\{\begin{array}{l}
            \partial_t x(a,t)+\partial_a x(a,t)=-\mu(a)x(a,t) \\
            x(0,t) = \displaystyle{\int_0^{a_{\max}}}\beta(a)x(a,t)da
            \end{array}\right.
        \end{equation*}
    \end{itemize}
\end{frame}

\footnotesize\tableofcontents

\section{Le modèle et sa dérivation}
% =====================================
% Le modèle et sa dérivation
% =====================================

\subsection{Modèle de Malthus}
\begin{frame}
    \begin{itemize}
        \item On considère une population homogène, isolée, qui vit dans un environnement stable avec des ressources illimitées. On note $P(t)$ le nombre total d'individu à l'instant $t \in \bbR_+$. \pause
        %\item Dans ces conditions, les variations de la population vont uniquement dépendre des taux de fertilité et de mortalité, notés respectivement $\beta$ et $\mu$, et supposés constants. \pause
        \item On observe la variation de population sur un intervalle de la forme $[t,t+h]$. On a 
        \begin{equation*}
            P(t+h) = P(t) + \text{naissances} - \text{décès}
        \end{equation*} \pause
        \item Sur l'intervalle $[t,t+h]$, le nombre de naissance est $\beta P(t)h$ et le nombre de décès est $\mu P(t) h$. On substitue dans l'équation précédente et on obtient
        \begin{equation*}
            P(t+h)=P(t)+\beta P(t)h - \mu P(t)h
        \end{equation*} \pause
        \item En divisant par $h$ et en faisant tendre $h$ vers $0$, on obtient finalement \textcolor{astral}{\textit{l'équation de Malthus}}
        \begin{equation*}
            P'(t) = (\beta - \mu)P(t)
        \end{equation*} %\pause
        %\item La solution de cet EDO est donnée par $P(t) = P(0)\ex^{(\beta - \mu)t}$.
    \end{itemize}
\end{frame}

\subsection{Modèle de McKendrick-Lotka}
\begin{frame}
    \begin{itemize}
        \item On va décrire l'évolution de la population par une fonction donnant la \textcolor{astral}{\textit{densité d'individu d'âge $a$ à l'instant $t$}}
        \begin{equation*}
            x(a,t) \geq 0, \ \ \ \ \ a \in [0,a_{\max}], \ t \geq 0
        \end{equation*}
        %où $a_{\max}$ est l'âge maximal de la population, supposé fini. 
        \pause
        \item On introduit la fonction $\beta$ qui représente le \textcolor{astral}{\textit{taux de fertilité de la population}} en fonction de l'âge. On peut alors considérer le nombre total de nouveau né à l'instant $t$
        \begin{equation*}
            B(t):= \int_{0}^{a_{\max}}\beta(a)x(a,t)d a
        \end{equation*} \pause
        \item On introduit à présent la fonction $\mu$ qui représente le \textcolor{astral}{\textit{taux de mortalité}} de la population en fonction de l'âge. On définit ensuite la \textcolor{astral}{\textit{probabilité de survie}}
        \begin{equation*}
            \Pi (a) := \exp\bigg(-\int_{0}^a \mu(\sigma)d \sigma\bigg)
        \end{equation*}
        et on pose $\Pi(a_{\max}) = 0$. \pause
        \item On définit de plus la \textcolor{astral}{\textit{fonction de maternité}} pour tout $a \in [0,a_{\max}]$, $K(a) := \beta(a)\Pi(a)$. \pause 
        \item Enfin, on définit le \textcolor{astral}{\textit{nombre de reproduction}}
        \begin{equation*}
            R := \int_0^{a_{\max}} \beta(a)\Pi(a)d a = \int_0^{a_{\max}} K(a)d a
        \end{equation*}
    \end{itemize}
\end{frame}

\begin{frame}
    \begin{itemize}
        \item On peut à présent définir le modèle de McKendrick-Lotka. Comme pour le modèle de Malthus, on étudie la variation de la population sur $[t,t+h]$
        \begin{equation*}
            x(a+h,t+h) = x(a,t) - \text{décès}
        \end{equation*} \pause 
        \item Le taux de décès entre $t$ et $t+h$ est donné par
        \begin{equation*}
            \int_0^{h}\mu(a+s)x(a+s,t+s)ds = \int_t^{t+h}\mu(a+s-t)x(a+s-t,s)d s
        \end{equation*} \pause  
        %\item On pose $G(t) := \displaystyle{\int_0^t}\mu(a+s-t)x(a+s-t,s)ds$, d'où
        %\begin{equation*}
        %    x(a+h,t+h)-x(a,t+h)+\big(x(a,t+h) - x(a,t)\big) = - \big(G(t+h)-G(t)\big)
        %\end{equation*} \pause
        \item On divise par $h$ et on fait tendre $h$ vers $0$, on obtient alors \textcolor{astral}{\textit{l'équation de McKendrick-Lotka}}
        \begin{equation*}
            \partial_tx(a,t)+\partial_ax(a,t) = -\mu(a)x(a,t)
        \end{equation*} \pause
        \item On obtient le \textcolor{astral}{\textit{modèle de McKendrick-Lotka}}
        \begin{equation*}
            \left\{\begin{array}{l}
            \partial_t x(a,t) + \partial_a x(a,t) = -\mu(a)x(a,t) \\
            x(0,t) = B(t) = \displaystyle{\int_0^{a_{\max}}} \beta (a)x(a,t)d a \\
            x(a,0) = x_0(a)
            \end{array}\right.
        \end{equation*}
    \end{itemize}
\end{frame}

\section{Résolution par la méthode des caractéristiques}
% ==================================================
% Résolution par la méthode des caracteristiques
% ==================================================

\subsection{Méthode des caractéristiques}
\begin{frame}
    \begin{itemize}
        \item Soit $a(X,t)$ la courbe caractéristique issue de $X \in \bbR$ fixé. La variation temporelle de $x$ le long de cette courbe s'écrit 
        \begin{eqnarray*}
            \frac{d}{d t}x\big(a(X,t),t\big) = \frac{\partial a}{\partial t}(X,t)\frac{\partial x}{\partial a}\big(a(X,t),t\big)+\frac{\partial x}{\partial t}\big(a(X,t),t \big)
        \end{eqnarray*} \pause
        \item En posant $\dfrac{\partial a}{\partial t}(X,t) \equiv 1$ on obtient 
        \begin{eqnarray*}
            \frac{d}{d t}x\big(a(X,t),t\big) = \frac{\partial x}{\partial a}\big(a(X,t),t\big)+\frac{\partial x}{\partial t}\big(a(X,t),t \big) = -\mu\big(a(X,t)\big)x\big(a(X,t),t\big)
        \end{eqnarray*} \pause
        \item Déterminons $a$, on a 
        \begin{eqnarray*}
            \frac{\partial a}{\partial t}(X,t) \equiv 1 &\Longrightarrow& a(X,t) = t + a(X,0)
        \end{eqnarray*}
    \end{itemize}
\end{frame}

\begin{frame}
    \begin{figure}[!ht]
        \centering
        \begin{subfigure}{0.45\textwidth}
            \centering
            \includegraphics[width=\textwidth]{images/caracteristiques/caracteristiques1.png}
            \caption{Avec condition au bord}
        \end{subfigure}
        \hfill
        \begin{subfigure}{0.45\textwidth}
            \centering
            \includegraphics[width=\textwidth]{images/caracteristiques/caracteristiques2.png}
            \caption{Sans condition au bord}
        \end{subfigure}
        \caption{Caractéristiques dans le demi-plan $(a,t)$ avec et sans condition au bord}
        \label{fig:caracteristiques}
    \end{figure} \pause
    \begin{itemize}
        \item Finalement, pour tout $X \in \bbR$
        \begin{eqnarray*}
            a(X,t) = t+X &\Longrightarrow& X = a-t
        \end{eqnarray*}
        et on a les conditions $t \in \bbR_+$ si $X \geq 0$ et $t \in [-X,+\infty[$ si $X<0$.
    \end{itemize}
\end{frame}

\begin{frame}
    \begin{itemize}
        \item La solution générale est
        \begin{equation*}
            x\big(a(X,t),t\big) = x\big(a(X,t_0),t_0\big)\exp\bigg(-\int_{t_0}^t\mu\big(a(X,s)\big)d s\bigg)
        \end{equation*}
        où $t_0 \in \bbR$ est le temps de départ de la caracteristique. \pause
        \item Si $t_0=0$, on a $X \geq 0$ donc $a \geq t$. La solution est alors donnée par 
        \begin{eqnarray*}
            x(a,t) = x_0(a-t)\exp\bigg(-\int_{0}^t\mu(a-t+s)d s\bigg)
        \end{eqnarray*} \pause
        \item Si $t_0=-X$, on a $X<0$ et donc $a<t$. La solution est alors donnée par 
        \begin{eqnarray*}
            x(a,t) = x(0,t-a)\exp\bigg(-\int_{t-a}^t\mu(a-t+s)d s\bigg)
        \end{eqnarray*} \pause
        \item La solution obtenue par la méthode des caractéristiques est donc
        \begin{equation*}
            x(a,t) = \left\{\begin{array}{lc}
                x_0(a - t)\dfrac{\Pi(a)}{\Pi(a-t)} & \text{ si } a \geq t, \\
                & \\
                B(t-a)\Pi(a) & \text{ si } a < t.
            \end{array}\right.
        \end{equation*}
    \end{itemize}
\end{frame}

\subsection{L'équation de renouvellement en démographie}
\begin{frame}
    \begin{itemize}
        \item On peut réinjecter son expression dans la condition au bord. On obtient
        \begin{eqnarray*}
            B(t) &=& \int_0^{a_{\max}} \beta(a)x(a,t)d a \\
            &=& \int_0^t\beta(a)B(t-a)\Pi(a)d a + \int_t^{a_{\max}} \beta(a)x_0(a-t)\frac{\Pi(a)}{\Pi(a-t)}d a \\
            &=& \int_0^t K(a)B(t-a)d a + \underbrace{\int_0^{a_{\max}} \beta(a+t)x_0(a)\frac{\Pi(a+t)}{\Pi(a)}d a}_{=:F(t)}
        \end{eqnarray*} \pause
        \item Finalement, on obtient \textcolor{astral}{\textit{l'équation de renouvellement}}
        \begin{equation*}
            B(t) = F(t)+\int_0^tK(a)B(t-a)d a
        \end{equation*}
    \end{itemize}
\end{frame}

\begin{frame}
    \begin{theoreme}
        L'équation de renouvellement admet une unique solution continue $B$ qui satisfait 
        \begin{equation*}
            \forall t \in \bbR_+, \ 0 \leq B(t) \leq \beta_+\ex^{\beta_+ t}\|x_0\|_{L^1}
        \end{equation*}
    \end{theoreme}
    \pause
    \begin{theoreme}
        Si $B$ est solution de l'équation de renouvellement alors $x$ définie par la méthode des caractéristiques est l'unique solution du modèle de McKendrick-Lotka. De plus, $x$ est continue sur $[0,a_{\max}] \times \bbR_+$ avec $a<t$. Enfin, pour $(a,t) \in [0,a_{\max}] \times \bbR_+$ on a $x(a,t) \geq 0$,
        \begin{equation*}
            x \in \scrC^0\big([0,T], L^1(0,a_{\max})\big)
        \end{equation*}
        et l'estimation
        \begin{equation*}
            \|x(t)\|_{L^1} \leq \ex^{\beta_+t}\|x_0\|_{L^1}
        \end{equation*}
    \end{theoreme}
\end{frame}

\subsection{Comportement asymptotique}
\begin{frame}
    \begin{itemize}
        \item On va réécrire l'équation de renouvellement sous forme de convolution
        \begin{equation*}
            B(t) = F(t) + (K*B)(t)
        \end{equation*} \pause
        \item Puis on applique la transformée de Laplace à celle-ci, ce qui donne 
        \begin{equation*}
            \hat{B}(z) = \hat{F}(z) + \hat{K}(z)\hat{B}(z)
        \end{equation*} \pause
        \item Par les propriétés de la transformée de Laplace, $\hat{K}$ et $\hat{F}$ sont analytiques donc les propriétés de $\hat{B}$ découlent du nombre de zéros de $1-\hat{K}$, autrement dit il nous faut résoudre l'équation suivante dans $\bbC$
        \begin{equation*}
            \hat{K}(z) = 1
        \end{equation*} \pause
        \item On considère la fonction réelle
        \begin{equation*}
            \hat{K}(x) = \int_0^{+\infty} \ex^{-xs}K(s)d s
        \end{equation*}
    \end{itemize}
\end{frame}

\begin{frame}
    \begin{itemize}
        \item La fonction réelle $\hat{K}$ est strictement décroissante et vérifie 
        \begin{eqnarray*}
            \lim_{x \to -\infty} \hat{K}(x) = +\infty \ \ \ &\text{ et }& \ \ \ \lim_{x \to +\infty} \hat{K}(x) = 0
        \end{eqnarray*}
        Par le TVI, il existe un unique $\alpha^* \in \bbR$ tel que $\hat{K}(\alpha^*) = 1$. \pause
        \item Et par stricte décroissance de $\hat{K}$ que 
        \begin{eqnarray*}
            \alpha^* >0 &\iff& \hat{K}(\alpha^*) = 1 < \hat{K}(0) = R \\
            \alpha^* < 0 &\iff& \hat{K}(\alpha^*) = 1 > \hat{K}(0)=R
        \end{eqnarray*} \pause
        \item Toute solution de l'équation $\hat{K}(z)=1$ admet une partie réelle strictement plus petite que $\alpha^*$. \pause
        \item On a $B(t) = B_0\ex^{\alpha^*t}\big(1+\Omega(t)\big)$ avec $B_0 \geq 0$ et $\displaystyle{\lim_{t \to +\infty}} \Omega(t)=0$. \pause On injecte dans la solution du modèle de McKendrick-Lotka et on obtient 
        \begin{equation*}
            \lim_{t \to +\infty} x(a,t) = \lim_{t \to +\infty} B_0\ex^{\alpha^*(t-a)}\Pi(a)
        \end{equation*}
    \end{itemize}
\end{frame}

%\begin{frame}[label={sec:org9c62e72}]{}
%\begin{block}{Titre et texte encadré}
%Du texte et encore du texte\ldots{} 
%\end{block}
%\end{frame}

\section{Résolution avec les semi-groupes}
% ======================================
% Résolution avec les semi-groupes
% ======================================

\subsection{La théorie des semi-groupes}
\begin{frame}
    Soit $X$ un espace de Banach.
    \begin{definition}
        Soit $\big\{T(t)\big\}_{t \geq 0}$ une famille d'opérateurs bornés de $X$ dans $X$, parametrée par un réel positif $t$. On dit que la famille $\big\{T(t)\big\}_{t \geq 0}$ est un \textcolor{astral}{semi-groupe d'opérateurs bornés sur $X$} si 
        \begin{itemize}
            \item[$(i)$] $T(0)=I$,
            \item[$(ii)$] $T(t+s)=T(t)T(s)$ pour tout $t,s \in \bbR_+$ (propriété des semi-groupes). 
        \end{itemize}
    \end{definition}
    \pause
    \begin{definition}
        Soit $\big\{T(t)\big\}_{t \geq 0}$ un semi-groupe d'opérateurs bornés. L'opérateur $A$ défini par 
        \begin{equation*}
            D(A) = \Big\{ x \in X \ \ | \ \ \lim_{h \to 0} \frac{T(h)x-x}{h} \ \text{ existe } \Big\}
        \end{equation*}
        et pour $x \in D(A)$
        \begin{equation*}
            Ax = \lim_{h \to 0} \frac{T(h)x-x}{h}
        \end{equation*}
        est le \textcolor{astral}{générateur infinitésimal du semi-groupe $\big\{T(t)\big\}_{t \geq 0}$}, et $D(A)$ est le domaine de définition de $A$.
    \end{definition}
\end{frame}

\begin{frame}
    \begin{definition}
        Soit $\big\{T(t)\big\}_{t \geq 0}$ un semi-groupe d'opérateurs bornés. On dit que le semi-groupe $\big\{T(t)\big\}_{t \geq 0}$ est \textcolor{astral}{continu au sens fort} si 
        \begin{equation*}
            \forall x \in X, \ \lim_{t \to 0} T(t)x=x
        \end{equation*}
        On dit que le semi-groupe $\big\{T(t)\big\}_{t \geq 0}$ est un \textcolor{astral}{$\scrC^0$ semi-groupe}.
    \end{definition}
    \pause
    \begin{definition}
        Soit $A$ un opérateur de $X$ dans $X$. 
        \begin{itemize}
            \item On appelle \textcolor{astral}{ensemble resolvant de $A$}, l'ensemble $\rho(A)$ des $\lambda \in \bbC$ tels que $\lambda I-A$ est inversible, \textit{i.e.} $(\lambda I-A)^{-1}$ est un opérateur borné de $X$.
            \item La famille $R(\lambda : A):=(\lambda I-A)^{-1}$, où $\lambda \in \rho(A)$, est appelée la \textcolor{astral}{résolvante de $A$}.
        \end{itemize}
    \end{definition}
\end{frame}

\subsection{Cadre théorique et problème de Cauchy}
\begin{frame}
    \begin{itemize}
        \item On va \og regarder \fg{} notre inconnue $[0,a_{\max}] \times \bbR_+ \ni (a,t) \mapsto x(a,t) \in \bbR_+$ comme une fonction $\bbR_+ \ni t \mapsto \tilde{x}(t)$ à valeurs dans un espace vectoriel de fonctions définies sur $\bbR_+$, c'est-à-dire $[0,a_{\max}] \ni a \mapsto \tilde{x}(t)(a) \in \bbR_+$. \pause
        \item On peut réécrire l'équation de McKendrick-Lotka et ainsi définir un opérateur $A$
        \begin{equation*}
            \frac{d}{d t}x(t) = \big(\underbrace{-\mu - \partial_a}_{=:A}\big) x(t)
        \end{equation*} \pause
        \item On obtient alors le problème de Cauchy
        \begin{equation*}
            \left\{\begin{array}{lr}
                \dfrac{d}{d t}x(t) = Ax(t) & \ \ \ t >0 \\
                x(0) = x_0 &
            \end{array}\right.
        \end{equation*} \pause
        \item Il nous faut intégrer la condition au bord du modèle dans le domaine de définition de $A$. Pour ça on pose
        \begin{equation*}
            D(A) := \Big\{ u \in L^1(0,a_{\max}) \ \ | \ \ u'+\mu u \in L^1(0,a_{\max}) \ \text{ et } \ u(0) = \int_0^{+\infty}\beta(a)u(a)\d a \Big\}
        \end{equation*}
    \end{itemize}
\end{frame}

\subsection{Application du théorème de Hille-Yosida}
\begin{frame}
    \begin{proposition}
        $A : D(A) \rightarrow L^1(0,a_{\max})$ est un opérateur fermé.
    \end{proposition}
    \pause
    \begin{proposition}
        $D(A)$ est dense dans $L^1(0,a_{\max})$.
    \end{proposition}
    \pause
    \begin{proposition}
        On a $]\|\beta\|_{L^{\infty}},+\infty[ \ \subset \rho(A)$ et pour chaque $\lambda$ dans cet ensemble la résolvante de $A$ vérifie
        \begin{equation*}
            \vertiii{R(\lambda : A)}_{L^1} \leq \frac{1}{\lambda - \|\beta\|_{L^\infty}}
        \end{equation*}
        où $\vertiii{\cdot}_{L^1}$ désigne la norme d'opérateur de $L^1$.
    \end{proposition}
    \pause
    \begin{itemize}
        \item Les trois propositions précédentes nous assurent que l'opérateur $A$ vérifie les hypothèses du théorème de Hille-Yosida, donc $A$ est le générateur d'un $\scrC^0$ semi-groupe que l'on note $\big\{T(t)\big\}_{t \geq 0}$ vérifiant 
        \begin{equation*}
            \forall t \in \bbR_+, \ \vertiii{T(t)}_{L^1} \leq \exp\big({\|\beta\|_{L^\infty}t}\big)
        \end{equation*}
    \end{itemize}
\end{frame}

\begin{frame}
    \begin{theorem}
        Le problème de Cauchy de McKendrick-Lotka admet une unique solution $x \in \scrC^0\big(\bbR_+, L^1(0,a_{\max})\big)$.
    \end{theorem}
    \pause
    \begin{itemize}
        \item \textbf{Existence et unicité dans $D(A)$.}\pause
        \begin{itemize}
            \item[$\triangleright$] Soit $x_0 \in D(A)$. On pose alors pour tout $t \in \bbR_+$
            \begin{equation*}
                \forall t \in \bbR_+, \ x(t) = T(t)x_0 \in D(A)
            \end{equation*}\pause
            \item[$\triangleright$] Alors avec les propriétés des semi-groupes, on obtient 
            \begin{eqnarray*}
                \frac{d}{d t}x(t) = AT(t)x_0 = Ax(t)
            \end{eqnarray*}\pause
            \item[$\triangleright$] On considère une autre solution $v$ du problème de Cauchy et on pose pour $t \in \bbR_+$ fixé
            \begin{equation*}
                \forall s \in \bbR_+, \ w(s) = T(t-s)v(s)
            \end{equation*}
        \end{itemize}
    \end{itemize}
\end{frame}

\subsection{Identification du semi-groupe}
\begin{frame}
    \begin{itemize}
        \item On a l'expression de la résolvante. Pour tout $\mu \in \ ]\|\beta\|_{L^\infty}, +\infty[$ et 
        \begin{equation*}
            R(\mu : A)x = \int_0^{+\infty}\ex^{-\mu t}T(t)x\d t
        \end{equation*} \pause
        \item En remplaçant par l'expression de la résolvante dans l'équation précédente et en identifiant, on obtient
        \begin{equation*}
            \forall t \in \bbR_+, \ \forall f \in L^1(0,a_{\max}), \ \big(T(t)f\big)(a) := \left\{\begin{array}{cl}\dfrac{\Pi(a)}{\Pi(a-t)}f(a-t) & \text{si } t \leq a, \\
            & \\
            \Pi(a)B(t-a) & \text{si } t > a.
            \end{array}\right.
        \end{equation*} \pause 
        \item La dépendance en $f$ dans l'expression de $T(t)$ dans le cas $t>a$ est implicite.
    \end{itemize}
\end{frame}

\subsection{Comportement asymptotique}
\begin{frame}
    \begin{definition}
        On définit la \textcolor{astral}{borne spectrale de $A$} par 
        \begin{equation*}
            s(A):= \sup\{\Re (\lambda) \in \bbR \ \ | \ \ \lambda \in \sigma(A)\}
        \end{equation*}
        avec $s(A)=-\infty$ si le spectre de $A$ est vide.
    \end{definition}
    \pause
    \begin{definition}
        On dit que le semi-groupe $\big\{T(t)\big\}_{t \geq 0}$ est \textcolor{astral}{irréductible} si $\{0\}$ et $X$ sont les seuls idéaux fermés $\big\{T(t)\big\}_{t \geq 0}$ invariant.
    \end{definition}
    \pause
    \begin{theorem}
        Soit $\big\{T(t)\big\}_{t \geq 0}$ un $\scrC^0$ semi-groupe positif tel que $\omega_{\text{ess}} < \omega_0$. Soit $\varepsilon >0$ donné par le théorème précédent. Si $\big\{T(t)\big\}_{t \geq 0}$ est irréductible et on a $x \in X^+$ et $x^* \in (X')^+$ tels que 
        \begin{equation*}
            \begin{array}{ccccc}
                Ax = s(A)x && A^*x^* = s(A)x^* && \langle x,x^*\rangle =1
            \end{array}
        \end{equation*}
        alors pour tout $\eta \in \ ]0,\varepsilon[$ et pour tout $x_0 \in X$, il existe une constante $M_\eta \geq 1$ telle que 
        \begin{equation*}
            \|\ex^{-s(A)t}T(t)x - \langle x,x_0^*\rangle x_0\| \leq M_\eta \ex^{-\eta t}\|x\|
        \end{equation*}
    \end{theorem}
\end{frame}

\begin{frame} 
    \begin{itemize}
        \item $\omega_{\text{ess}}=-\infty$ et $\omega_0 = s(A) = \alpha^*$ \pause
        \item $x(a) = x(0)\ex^{-\alpha^*a}\Pi(a)$ \pause
        \item $A^*v = -\mu v + v' + \beta v(0)$ \pause
        \item $x^*(a) = x^*(0)\dfrac{\ex^{\alpha^*a}}{\Pi(a)}\Big(1-\displaystyle{\int_0^a}\ex^{-\alpha^*s}K(s)d s\Big)$ \pause 
        \item Pour satisfaire la condition $\langle x,x^*\rangle =1$, on peut par exemple prendre $x(0)=1$ et déterminer $x^*(0)$ satisfaisant l'équation. \pause
    \end{itemize}
    \begin{proposition}
        Le semi-groupe $\big\{T(t)\big\}_{t \geq 0}$ est irréductible si et seulement si 
        \begin{equation*}
            \sup \ ]0,a_{\max}[ \ \setminus \omega = a_{\max}
        \end{equation*}
        où $\omega \subset \ ]0,a_{\max}[$ est l'ouvert constitué des points au voisinage desquels $\beta=0$.
    \end{proposition}
    \pause
    Donc on a par le théorème
    \begin{equation*}
        \lim_{t \to +\infty} \big(T(t)x_0\big)(a) = \lim_{t \to +\infty} \ex^{\alpha^*(t-a)}\Pi(a)\langle x_0,x^*\rangle
    \end{equation*}
\end{frame}

%\begin{frame}[label={sec:org88fd8be}]{Avec un titre plus imposant}
%Du texte non encadré \pause 
%\begin{block}{Un sous-titre}
%Du texte et une référence bibliographique au sujet de, par exemple, Quine.
%\end{block}
%\end{frame}

\section{Modèle non-linéaire}
% ==========================
% Modèle non-linéaire
% ==========================

\subsection{Modèle de Verhulst}
\begin{frame}
    \begin{itemize}
        \item On commence par étudier le modèle EDO introduit par Pierre François Verhulst qui généralise le modèle de Malthus. \pause
        \item Le modèle de Verhulst fait l'hypothèse que le taux de natalité et le taux de mortalité sont des fonctions affines respectivement décroissante et croissante de la taille de la population. \pause
        \item En s'inspirant du modèle de Malthus, on obtient l'équation 
        \begin{equation*}
            P'(t) = \big(\beta(P)-\mu(P)\big)P(t)
        \end{equation*} \pause
        \item On pose $\beta(P)-\mu(P)=a-bP$, avec $a,b>0$. L'équation devient 
        \begin{equation*}
            P'(t) = \big(a-bP(t)\big)P(t) = aP(t)-bP(t)^2
        \end{equation*} \pause
        \item En posant $K=a/b$ la \textcolor{astral}{\textit{capacité d'accueil}}, on obtient la forme \og logistique \fg{} de l'équation 
        \begin{equation*}
            P'(t) = aP(t)\Big(1-\frac{P(t)}{K}\Big)
        \end{equation*}
    \end{itemize}
\end{frame}

\begin{frame}
    \begin{itemize}
        \item La solution du modèle de Verhulst est donnée par 
        \begin{equation*}
            P(t) = K \frac{1}{1+\Big(\dfrac{K}{P(0)}-1\Big)\ex^{-at}}
        \end{equation*} \pause
        \item On peut montrer que le modèle de Verhulst admet deux points d'équilibres : $0$ et $K$. \pause
        \item On étudie la stabilité de ces points d'équilibres. \pause
        \item $0$ est stable mais n'est pas localement asymptotiquement stable. \pause
        \item $K$ est stable et localement asymptotiquement stable.
    \end{itemize}
\end{frame}

\subsection{Modèle de McKendrick-Lotka non-linéaire}
\begin{frame}
    \begin{itemize}
        \item On pourrait considérer l'équation 
        \begin{equation*}
            \partial_tx(a,t)+\partial_ax(a,t) = -\mu(a)x(a,t)-\gamma(a)x(a,t)^2
        \end{equation*}
        où $\gamma \in L^\infty(0,a_{\max})$ et $\gamma > 0$, cette fonction peut représenter un terme de compétition inter-population (pour les ressources par exemple). \pause
        \item On va plutôt considérer l'équation 
        \begin{equation*}
            \partial_t x(a,t) + \partial_a x(a,t) = -\mu(a)x(a,t)-\gamma(a)\|x(\cdot,t)\|_{L^1}x(a,t)
        \end{equation*} \pause
        \item On considère alors le modèle non-linéaire de McKendrick-Lotka suivant
        \begin{equation*}
            \left\{\begin{array}{l}
                \partial_t x(a,t) + \partial_a x(a,t) = -\mu(a)x(a,t)-\gamma(a)\|x(\cdot,t)\|_{L^1}x(a,t) \\
                x(0,t) = B(t) = \displaystyle{\int_0^{a_{\max}}} \beta (a)x(a,t)d a \\
                x(a,0) = x_0(a)
            \end{array}\right.
        \end{equation*}
    \end{itemize}
\end{frame}

\begin{frame}
    \begin{itemize}
        \item Comme dans la section précédente, on considère $x: \bbR_+ \rightarrow L^1(0,a_{\max})$ et l'opérateur $A$. Et donc on a 
        \begin{equation*}
            \left\{\begin{array}{l}
                \dfrac{d}{d t}x(t) = Ax(t) -\gamma\|x(t)\|_{L^1}x(t) \\
                x(0) = x_0
            \end{array}\right.
        \end{equation*}\pause
        \item On pose alors 
        \begin{eqnarray*}
            f : L^1(0,a_{\max}) \rightarrow L^1(0,a_{\max}), \ u \mapsto - \gamma \|u\|_{L^1}u
        \end{eqnarray*}
        de sorte à obtenir le problème 
        \begin{equation*}
            \left\{\begin{array}{l}
                \dfrac{d}{d t}x(t) = Ax(t) + f\big(t,x(t)\big) \\
                x(0) = x_0
            \end{array}\right.
        \end{equation*}
    \end{itemize}
\end{frame}

\begin{frame}
    \begin{definition}
        Une solution continue $u$ de l'équation intégrale
        \begin{equation*}
            u(t) = T(t-t_0)u_0 + \int_{t_0}^t T(t-s)f\big(s,u(s)\big)d s
        \end{equation*}
        est appelée \textcolor{astral}{solution mild}.
    \end{definition}
    \pause
    \begin{theorem}
        Soit $f:[t_0,T] \times X \rightarrow X$ continue en $t$ sur $[t_0,T]$ et uniformément Lipschitz de constante $L$ sur $X$. Si $-A$ est le générateur infinitésimal d'un $\scrC^0$ semi-groupe $\big\{T(t)\big\}_{t \geq 0}$ alors pour toute condition initiale $u_0 \in X$, le problème
        \begin{equation*}
            \left\{\begin{array}{l}
                \dfrac{d}{d t}u(t) +Au(t)=f\big(t,u(t)\big) \ \ \ \ \ t >t_0 \\
                u(t_0) = u_0
            \end{array}\right.
        \end{equation*}
        admet une unique solution mild $u \in \scrC^0\big([t_0,T],X\big)$. De plus, l'application $u_0 \mapsto u$ est Lipschitzienne de $X$ dans $\scrC^0\big([t_0,T],X\big)$.
    \end{theorem}
    \pause
    \begin{theorem}
        Le problème de McKendrick-Lotka non-lineaire admet une unique solution mild dans $L^1(0,a_{\max})$.
    \end{theorem}
\end{frame}

\begin{frame}{Points d'équilibres}
    \begin{itemize}
        \item Déterminer les points d'équilibres revient à chercher une solution $x_0 \in D(A)$ de notre EDP qui ne dépend pas du temps. \pause
        \item En injectant $x_0$ dans l'EDP non-linéaire, on obtient 
        \begin{equation*}
            Ax_0+f(x_0) = 0
        \end{equation*} \pause
        \item La solution de cette EDO est donnée par 
        \begin{equation*}
            x_0(a) = x_0(0)\Pi(a)\exp\bigg(-\|x_0\|_{L^1}\int_0^a\gamma(\sigma)d \sigma\bigg)
        \end{equation*} \pause
        \item Comme $x_0 \in D(A)$, en injectant l'expression de $x_0$ dans la condition au bord, on obtient
        \begin{equation*}
            x_0(0)\bigg(1-\int_0^{a_{\max}}K(a)\exp\bigg(-\|x_0\|_{L^1}\int_0^a \gamma(\sigma)d \sigma\bigg)d a\bigg) = 0
        \end{equation*}
    \end{itemize}
\end{frame}

\begin{frame}{Points d'équilibres}
    \begin{itemize}
        \item Si $x_0(0)=0$, alors $x_0\equiv 0$ et donc le seul point d'équilibre est $0$. \pause
        \item Sinon, on a 
        \begin{equation*}
            \int_0^{a_{\max}}K(a)\exp\bigg(-\|x_0\|_{L^1}\int_0^a\gamma(\sigma)\d \sigma\bigg)\d a = 1
        \end{equation*} \pause
        \item On étudie la fonction 
        \begin{equation*}
            g : \bbR_+ \rightarrow \bbR, \ y \mapsto \int_0^{a_{\max}}K(a)\exp\bigg(-y\int_0^a\gamma(\sigma)\d \sigma\bigg)\d a
        \end{equation*} \pause
        \item $g$ est positive, continue, strictement décroissante et vérifie 
        \begin{eqnarray*}
            g(0) = \int_0^{a_{\max}}K(a)\d a = R \ \ \ &\text{ et }& \ \ \ \lim_{y \to +\infty}g(y) = 0
        \end{eqnarray*} \pause
        \item Si $R\leq1$, alors il n'y a pas de point d'équilibre non-trivial. Si $R > 1$ alors, on a par le TVI l'existence d'une unique solution $y^*\geq 0$ à l'équation et donc d'un unique point d'équilibre vérifiant
        $\|x_0\|_{L^1}=y^*$.
    \end{itemize}
\end{frame}

\begin{frame}
    \begin{definition}
        On dit que $x^* \in X$ est un \textcolor{astral}{point d'équilibre localement exponentiellement asymptotiquement stable} s'il existe $\varepsilon >0$, $M \geq 1$ et $\omega <0$ tels que si $x \in X$ et $\|x-x^*\| \leq \varepsilon$ alors $t_{\max}=+\infty$ et 
        \begin{equation*}
            \forall t \geq 0, \ \|u(t,x)-x^*\| \leq M\ex^{\omega t}\|x-x^*\|
        \end{equation*}
    \end{definition}
    \pause
    \begin{proposition}
        Soit $\big\{T(t)\big\}_{t \geq 0}$ un $\scrC^0$ semi-groupe d'opérateurs sur $X$ de générateur infinitésimal $A$. Soit $f$ un opérateur non-linéaire de $X$ dans $X$ différentiable sur $X$. Pour chaque $x \in X$, on note $u(t,x)$ la solution mild du problème de Cauchy non-linéaire. Soit $x^*$ tel que $Ax^*+f(x^*)=0$ et soit $\big\{S(t)\big\}_{t\geq 0}$ le $\scrC^0$ semi-groupe d'opérateurs sur $X$ de générateur infinitésimal $\hat{A}:=A+f'(x^*)$. Alors on a 
        \begin{itemize}
            \item Si $\max\{\omega_{\text{ess}}(\hat{A}), \displaystyle{\sup_{\lambda \in \sigma(\hat{A})\setminus \sigma_{\text{ess}}(\hat{A})}}\Re(\lambda)\}<0$, alors $x^*$ est un point d'équilibre localement exponentiellement asymptotiquement stable.
            \item S'il existe $\lambda_1 \in \sigma(\hat{A})$ tel que $\Re(\lambda_1) >0$ et $\max\{\omega_{\text{ess}}(\hat{A}), \displaystyle{\sup_{\lambda \in \sigma(\hat{A})\setminus \sigma_{\text{ess}}(\hat{A}), \ \lambda \neq \lambda_1}}\Re(\lambda)\}<\Re(\lambda_1)$, alors $x^*$ est un point d'équilibre instable.
        \end{itemize}
    \end{proposition}
\end{frame}

\begin{frame}{Stabilité}
    \begin{itemize}
        \item On considère à présent l'opérateur $\hat{A} := A + f'(x_0)$.
    \end{itemize}
    \begin{proposition}
        L'opérateur $\hat{A}$ est le générateur d'un $\scrC^0$ semi-groupe.
    \end{proposition}
    \pause
    \begin{itemize}
        \item Comme dans le cas linéaire, on a $\omega_{\text{ess}} = -\infty$ et $\sigma_{\text{ess}}(\hat{A})=\emptyset$. \pause
        \item On pose 
        \begin{equation*}
            \Phi(a) := \exp\bigg(-\|x_0\|_{L^1}\int_0^a\gamma(\sigma)d \sigma\bigg)
        \end{equation*}
        cette fonction est une mortalité supplémentaire due à la compétition intra-population. \pause
    \end{itemize}
    \begin{theorem}
        On a 
        \begin{equation*}
            \sigma(\hat{A}) = \mathrm{VP}(\hat{A}) = \big\{\lambda \in \bbC \ | \ \widehat{K\Phi}(\lambda)=1\big\}
        \end{equation*}
        L'équation $\widehat{K\Phi}(\lambda)=1$ admet une unique solution réelle $\lambda^*$ qui vérifie $\lambda^*=s(\hat{A})$, c'est-à-dire que c'est la plus grande valeur propre de $\hat{A}$. De plus, on a $s(\hat{A}) \leq s(A)$.
    \end{theorem}
\end{frame}

\begin{frame}{Stabilité}
    \begin{itemize}
        \item On rappelle que, dans l'étude asymptotique du modèle linéaire, on a obtenu les équivalences suivantes 
        \begin{align*}
            \alpha^* < 0 \iff R < 1 & \quad & \alpha^* > 0 \iff R > 1 & \quad & \alpha^*=0 \iff R=1
        \end{align*} \pause
        \item Si $R<1$, alors $0$ est le seul point d'équilibre. De plus, on a $\alpha^*<0$ et notre opérateur est $\hat{A}=A+f'(0)=A$, donc $s(A) <0$ et $0$ est localement exponentiellement asymptotiquement stable. \pause
        \item Si $R=1$, alors $0$ est le seul point d'équilibre. De plus, on a $\alpha^*=0$ et notre opérateur est $\hat{A}=A+f'(0)=A$, donc $s(\hat{A})=s(A) =0$. On ne peut donc pas conclure sur la stabilité du point d'équilibre. \pause
    \end{itemize}
\end{frame}

\begin{frame}{Stabilité}
    \begin{itemize}
        \item Si $R>1$, alors on a deux points d'équilibre : $0$ et $x_0$. On a aussi $\alpha^*>0$ et notre opérateur est $\hat{A}=A+f'(0)=A$, donc $s(\hat{A})=s(A) >0$ par la proposition, donc $0$ est instable. \pause
        \item Pour $x_0$, qui est un point d'équilibre non-trivial, notre opérateur est $\hat{A}=A+f'(x_0)$ et $\alpha^*>0$. \pause
        \item Si $0 < \lambda^* \leq \alpha^*$ alors $x_0$ est instable. \pause
        \item Si $\lambda^*=0$ alors on ne peut pas conclure sur la stabilité de $x_0$. \pause
        \item Si $\lambda^* < 0$ alors $x_0$ est localement exponentiellement asymptotiquement stable.
    \end{itemize}
\end{frame}


%\section{Approximation par la méthode Galerkin discontinu}
%\begin{frame}
\begin{itemize}
    \item  Fait un lien entre les méthodes de volumes finis et d'éléments finis. \\
    Recherche de la formulation variationnelle puis on cherche la solution discrète.
    \item Cependant grande différence avec ces méthodes, aucune continuité n’est imposée aux interfaces entre les éléments du maillage.
\end{itemize}
   
\begin{figure}
    \centering
    \includegraphics[width=0.5\linewidth]{images/Galerkine_Discontinue/image_maille_GD.JPG}
    \caption{Exemple d'une façon de gérer le flux au niveau de l'intervalle.}
    \label{flux_GD_exemple}
\end{figure}
\begin{equation*}
    \{\!\!\{v\}\!\!\}_F(x) := \tfrac{1}{2}\left(v|_{T_1}(x)+v|_{T_2}(x)\right) \ \ \ \ \ \ \ \ [\![v]\!] := v|_{T_1}(x)-v|_{T_2}(x)
\end{equation*}

\end{frame}

\begin{frame}{Réécrire le problème}
\begin{itemize}
\item On commence par partir du problème de Rayleigh-Stokes $1$D homogène
\begin{equation*}
    \begin{cases}
        \partial_tu(t,x) - \nu\partial_{xx} u(t,x) - \alpha\partial_{txx} u(t,x) = 0 & (t,x) \in \mathbb R_+ \times \bbR  \\
        u(0,x) = u_0(x) & x \in \bbR
    \end{cases}
\end{equation*}  
\item Puis en faisant comme en élément fini, on se retrouve en $1D$ avec:
\begin{equation*}
    \int_{K_i} \partial_t uv + \nu \int_{K_i}\partial_{x}u\partial_{x}v \ + \alpha \int_{K_i} \partial_{tx}u\partial_{x}v
    =  \nu \Big[\partial_{x}u v \Big]^{x_i +h/2}_{x_i - h/2} \ + \alpha \Big[\partial_{tx}u v \Big]^{x_i +h/2}_{x_i - h/2} 
\end{equation*}
\item Le choix de la base change le problème. Prenons la base $\{e_0\} = \{1\}$.
\begin{equation*}
    \partial_x u_h\big|_{x_i + h/2} = \frac{u_0^{i+1} - u_0^i}{h}, \qquad
    \partial_x u_h\big|_{x_i - h/2} = \frac{u_0^i - u_0^{i-1}}{h}
\end{equation*}
\begin{equation*}
    h\,\partial_t u_0^i = \frac{\nu}{h}\left(u_0^{i+1} - 2u_0^i + u_0^{i-1}\right)
    + \frac{\alpha}{h}\left(\partial_t u_0^{i+1} - 2\partial_t u_0^i + \partial_t u_0^{i-1}\right)
\end{equation*}
\end{itemize}
\end{frame}

\begin{frame}{Le cas P1}
Pour $j$ valant $0$ ou $1$ on obtient le système suivant:
\begin{eqnarray*}\label{EDP_RS_Formulation_Faible_P1}
  &\partial_t u_i^0 \displaystyle{\int_{K_i}}  \phi_0 \phi_j + \partial_t u_i^1 \displaystyle{\int_{K_i}}   \phi_1 \phi_j& \\
   &+\nu \ u_i^0  \displaystyle{\int_{K_i}} \partial_{x}\phi_0 \partial_{x}\phi_j  + \nu \ u_i^1 \displaystyle{\int_{K_i}} \partial_{x}\phi_1 \partial_{x}\phi_j& \\
   &+ \alpha \partial_{t}u_i^0 \displaystyle{\int_{K_i}}  \partial_{x}\phi_0\partial_{x}\phi_j + \alpha \partial_{t}u_i^1 \displaystyle{\int_{K_i}}  \partial_{x}\phi_1\partial_{x}\phi_j &\\
    &=& \\ 
    &\nu \Big[\partial_{x}\big(u_i^0\phi_0 +u_i^1\phi_1 \big) \phi_j  \Big]^{x_i +h/2}_{x_i - h/2} \ + \alpha \Big[\partial_{tx}\big(u_i^0\phi_0 +u_i^1\phi_1 \big) \phi_j \Big]^{x_i +h/2}_{x_i - h/2} 
\end{eqnarray*}
Ce qui en calculant les matrices de rigidité et de masse donne:
\begin{eqnarray*}
    &&\begin{bmatrix}
        h & 0 \\
        0 & \frac{h}{12} 
    \end{bmatrix}
    \begin{bmatrix}
    \partial_t u_i^0 \\
    \partial_t u_i^1    
    \end{bmatrix}
    +
        \begin{bmatrix}
        0 & 0 \\
        0 & \frac{\nu + \alpha \partial_t}{h} 
    \end{bmatrix}
    \begin{bmatrix}
    u_i^0 \\
    u_i^1    
    \end{bmatrix}
    = (\nu + \alpha\partial_t)
    \begin{bmatrix}
    \{\!\!\{ u_i^1 \}\!\!\}_{F_i} - \{\!\!\{ u_i^1 \}\!\!\}_{F_{i-1}}\\
    \{\!\!\{ u_i^1 \}\!\!\}_{F_i} \times\frac{1}{2} - \{\!\!\{ u_i^1 \}\!\!\}_{F_{i-1}} \times (- \frac{1}{2})
    \end{bmatrix} \\
\end{eqnarray*}
\end{frame}

\begin{frame}
    \begin{itemize}
        \item Les problèmes de second grade doivent traiter de façon plus subtile le flux numérique.
        \item La méthode reste naturelle pour traiter ce genre de problème. Le terme $\alpha$ permet de régulariser les phénomènes de chocs qui entraînent une diffusion numérique pour cette méthode.
    \end{itemize}
\end{frame}

\begin{frame}
    \begin{center}
        \Large Merci pour votre écoute !
    \end{center}
\end{frame}

\end{document}